% =====================================================================
%  DATOS DEL TRABAJO  ·  rellena SOLO este fichero para personalizar
%  la memoria, los informes de seguimiento y la portada provisional.
% =====================================================================

% --- Tipo de trabajo (tal como aparece en la portada oficial, en catalán)
\newcommand{\tfgtipo}{Treball de Fi de Grau}       % o: Treball de Fi de Màster
\newcommand{\tfgsigla}{TFG}                        % o: TFM

% --- Estudiante y título (el título debe ser IDÉNTICO al de la propuesta
%     administrativa y al de la portada oficial)
\newcommand{\tfgautor}{Nombre Apellido1 Apellido2}
\newcommand{\tfgtitulo}{Título del trabajo: subtítulo que diga qué se ha hecho}

% --- Titulación y centro (en la lengua de la portada oficial: catalán)
\newcommand{\tfgtitulacion}{Grau en Enginyeria Informàtica}
\newcommand{\tfgcentro}{Escola Politècnica Superior}
\newcommand{\tfgcampus}{Campus d'Igualada}         % o: Campus de Cappont (Lleida)
\newcommand{\tfguniversidad}{Universitat de Lleida}

% --- Dirección. Deja \tfgcodirector vacío ({}) si no hay codirector/a.
\newcommand{\tfgdirector}{Nombre y apellidos del director/a}
\newcommand{\tfgcodirector}{}

% --- Presentación (mes en catalán, como en la portada oficial: Juny, Juliol, Setembre)
\newcommand{\tfgmes}{Juliol}
\newcommand{\tfganyo}{2027}
\newcommand{\tfgcurso}{2026--2027}

% --- Palabras clave (3-6), en los tres idiomas del resumen
\newcommand{\tfgclavesES}{palabra clave 1, palabra clave 2, palabra clave 3}
\newcommand{\tfgclavesCA}{paraula clau 1, paraula clau 2, paraula clau 3}
\newcommand{\tfgclavesEN}{keyword 1, keyword 2, keyword 3}

% --- Idioma principal del documento: spanish | catalan | english
%     (cambia "Capítulo", "Índice", "Figura", "Tabla"... automáticamente)
\newcommand{\idiomadocumento}{spanish}

% --- Estilo de citas: ieee (numérico [1], recomendado) | apa (autor-año)
%     Elige UNO y no lo mezcles nunca.
\newcommand{\estilocitas}{ieee}

% --- Enlaces que aparecen en el Apéndice A (repositorio y vídeo demo)
\newcommand{\tfgrepo}{https://github.com/usuario/repositorio}
\newcommand{\tfgversion}{v1.0.0}                   % tag o commit entregado
\newcommand{\tfgvideo}{https://youtu.be/XXXXXXXXXXX}
